\chapter{Cognitive Testing}
\section*{What Is This Test?}
Cognitive testing evaluates memory, attention, language, processing speed, reasoning, and visuospatial ability. It helps characterize changes and functional strengths; a single score does not identify one disease.
\section*{Alternative Names}
\begin{itemize}\item Neuropsychological Testing\item Cognitive Assessment\item Memory Testing\item Neurocognitive Evaluation\end{itemize}
\section*{Principle}
Standardized tasks are compared with expected performance for age, education, language, and other relevant factors. Brief screening tools identify concern; full neuropsychological batteries examine patterns across domains.
\section*{Why This Test Is Ordered}
It may evaluate memory complaints, suspected dementia or mild cognitive impairment, traumatic brain injury, stroke, epilepsy, developmental conditions, learning problems, or effects of psychiatric illness, sleep disorders, medicines, or substance use.
\section*{Primary vs Secondary Test}
History, examination, medication review, mood and sleep assessment, and functional history are primary. Cognitive testing is a secondary assessment that quantifies performance.
\section*{How This Test Is Performed}
A psychologist or trained clinician administers spoken, written, visual, and problem-solving tasks. Testing may take minutes for screening or several hours for a comprehensive evaluation.
\section*{What Preparation Is Needed}
Sleep normally, bring glasses or hearing aids, and provide education, language, medical history, medicines, and examples of the concern.
\section*{Understanding Results}
Scores are interpreted as a pattern. Anxiety, depression, fatigue, pain, language, culture, hearing, vision, and practice effects can influence performance. A low score alone does not establish dementia or incapacity.
\section*{Follow-up Tests}
Follow-up may include laboratory testing for reversible causes, brain MRI or CT, hearing and vision assessment, sleep evaluation, psychiatric assessment, or repeat testing.
\section*{References}
\begin{enumerate}\item MedlinePlus. Cognitive Testing.\item National Institute on Aging. Assessing cognitive impairment and dementia.\end{enumerate}
\clearpage\section*{Understanding Results and Follow-up}
The laboratory report should identify the specimen, method, units, reference interval or decision limit, and whether the result is positive, negative, abnormal, or inconclusive. Reference intervals vary with age, sex, pregnancy, specimen type, assay, and laboratory; a result outside a range is not automatically a diagnosis, and a result within a range does not always exclude disease. Discuss unexpected results with the ordering clinician, who can decide whether repeat or confirmatory testing, treatment, or referral is appropriate. Do not change medicines or delay urgent care based on an isolated result.

\section*{Regulatory Status and Authorization Date}This chapter describes a test category, not a specific commercial product. FDA, CDSCO, EU IVDR, and MHRA approval or registration dates therefore cannot be assigned without the assay manufacturer, product name, instrument, and intended use. For laboratory-developed tests, an individual agency approval date may not exist. See the Preface for the interpretation of regulatory status.
\clearpage
